% ============================================================
%  INTRODUCTION GÉNÉRALE
% ============================================================
\newpage
\addcontentsline{toc}{chapter}{Introduction générale}
\vspace*{1.5cm}

\begin{center}
    \textbf{\Large Introduction générale}
\end{center}
\vspace{1cm}

{
\setlength{\parindent}{0pt}
\setlength{\parskip}{1em}

\textbf{Contexte et enjeux}

La gestion des ressources en eau constitue l'un des défis majeurs du XXIe siècle, particulièrement dans les régions arides et semi-arides comme le Maroc. Face à la raréfaction des ressources hydriques et à l'augmentation des besoins liés à la croissance démographique et au développement agricole, la mobilisation et l'optimisation des ouvrages hydrauliques existants deviennent une priorité stratégique. Dans ce contexte, les barrages jouent un rôle fondamental pour la sécurisation de l'approvisionnement en eau potable, l'irrigation des terres agricoles et la production d'énergie hydroélectrique.

Le barrage Ahl Souss, également dénommé barrage Ait M'zal, constitue un ouvrage stratégique pour la région de Chtouka Ait Baha. Construit en 2005, cet ouvrage en béton compacté au rouleau (BCR) remplit des fonctions essentielles : alimentation en eau potable des populations, irrigation des périmètres agricoles et recharge de la nappe phréatique. Ses organes de sécurité, notamment l'évacuateur de crues et la vidange de fond, doivent garantir un fonctionnement fiable même lors des événements hydrologiques extrêmes.

\textbf{Problématique et objectifs}

La conception et la vérification du comportement hydraulique des ouvrages complexes comme l'évacuateur de crues à seuil latéral et la vidange de fond du barrage Ahl Souss nécessitent une analyse approfondie des écoulements. Les approches théoriques classiques, basées sur des formules empiriques et des hypothèses simplificatrices, montrent leurs limites face à la complexité géométrique et aux phénomènes physiques mis en jeu, notamment la turbulence, les écoulements tridimensionnels et les interactions fluide-structure.

La modélisation physique sur modèles réduits, bien que représentative de la réalité, présente des inconvénients majeurs : coûts élevés, délais de réalisation importants et effets d'échelle difficiles à maîtriser. Face à ces limitations, la mécanique des fluides numérique (CFD) s'impose comme une alternative complémentaire puissante, permettant une analyse détaillée des écoulements à moindre coût et dans des délais réduits.

Ce projet de fin d'études s'inscrit dans cette démarche et vise à :
\begin{itemize}
    \item \textbf{Modéliser numériquement} l'écoulement dans l'évacuateur de crues du barrage Ahl Souss.
    \item \textbf{Simuler le comportement hydraulique} de la vidange de fond.
    \item \textbf{Comparer les résultats numériques} avec les données théoriques et dimensionnelles.
    \item \textbf{Proposer des recommandations} pour l'optimisation des ouvrages étudiés.
\end{itemize}

\textbf{Méthodologie adoptée}

Pour atteindre ces objectifs, une approche méthodologique structurée a été adoptée :
\begin{enumerate}
    \item \textbf{Phase bibliographique} : analyse des fondements théoriques de la CFD, étude des modèles de turbulence, revue des méthodes de discrétisation et présentation détaillée du barrage Ahl Souss.
    \item \textbf{Phase d'analyse des outils} : benchmark des logiciels de simulation 3D et description approfondie d'OpenFOAM.
    \item \textbf{Phase de modélisation} : construction des géométries sous Autodesk Revit, génération des maillages, définition des conditions aux limites et paramétrage des solveurs.
    \item \textbf{Phase d'analyse des résultats} : interprétation des champs de vitesse et de pression, comparaison avec les données théoriques et formulation de recommandations.
\end{enumerate}

}
\newpage
